\section{Conclusão}
\label{sec:conclusion}

Este trabalho formalizou o Problema da Cobertura Mínima
por Vértices, apresentou a NP-Completude de sua versão
de decisão e comparou duas implementações em Java~17
puro: Branch and Bound e matching maximal
2-aproximado. Os testes independentes, os controles
estruturados e a preservação de sementes e medições
permitiram relacionar fundamentos e comportamento
experimental de maneira auditável.

Na execução analisada, o exato comprovou
\OptimalInstances{} ótimos entre \TotalInstances{}
instâncias e foi interrompido em duas; as três repetições
dessas instâncias explicam os seis timeouts individuais.
A aproximação apresentou menor tempo em todas as
comparações selecionadas. Nas \QualityInstances{}
instâncias com ótimo positivo comprovado, sua razão
média foi \MeanRatio{}, sem violação do limite~2.
Esse ganho de tempo veio acompanhado de coberturas
eventualmente maiores, enquanto a certificação do mínimo
impôs custos elevados ao exato em parte dos grafos.

Os objetivos foram atendidos no escopo da implementação
e do corpus utilizado, sem converter incumbentes após
timeout em ótimos comprovados. Permanecem limitações
de escala, de diversidade das instâncias e de controle
da JVM. Como trabalhos futuros, propõem-se melhores
incumbentes iniciais, limites inferiores mais fortes,
redução ou memorização de subproblemas, avaliação
separada das heurísticas, outras aproximações e estudo
de grafos reais. Múltiplas execuções independentes da
JVM e parametrizações pelo tamanho da cobertura
também permitiriam aprofundar a análise.
